\subsection{Sampling-Based Consistency Methods}

Self-consistency \cite{Wang2022SelfConsistency} introduced the idea of aggregating multiple chain-of-thought samples to improve reasoning accuracy, demonstrating that majority voting over stochastic samples can recover correct answers that individual samples miss. This work established the principle that generative variance can be harnessed as a signal, but focused on accuracy rather than hallucination detection.

Semantic uncertainty \cite{Kuhn2023Semantic} extended this to uncertainty quantification, proposing to cluster semantically equivalent samples and measure entropy over clusters as a proxy for model uncertainty. The key insight is that linguistic variation should be abstracted away --- two samples expressing the same proposition should count as agreeing --- making NLI-based semantic equivalence central to the approach. Kuhn et al.\ validate on open-ended QA with base language models.

SelfCheckGPT \cite{Manakul2023SelfCheckGPT} operationalized consistency-based hallucination detection for black-box LLMs by comparing each sample against the others using NLI, \texttt{BERTScore}, or $n$-gram overlap. Inconsistent samples are flagged as likely hallucinations. Validation is on WikiBio with GPT-3, a base-class model. Our work directly evaluates whether this paradigm transfers to instruction-tuned models.

\subsection{Uncertainty Quantification for LLMs}

\cite{Xiong2023Can} provide a comprehensive study of LLM calibration, showing that instruction-tuned models are systematically overconfident. Their findings suggest that RLHF fine-tuning may alter the relationship between expressed confidence and actual accuracy --- a precursor to the regime shift we characterize. \cite{Kadavath2022Language} study whether LLMs know what they know, finding that larger models are better calibrated but that calibration degrades on out-of-distribution inputs.

\cite{Lin2024Generating} study consistency-based uncertainty estimation under distribution shift, finding that consistency remains informative when the generative distribution is broad but degrades when the model is strongly conditioned. Our work identifies RLHF fine-tuning as one mechanism that triggers this degradation on structured factual tasks.

\cite{Huang2023Survey} survey the broader hallucination detection landscape, including sampling-based, retrieval-based, and probing-based methods. They note that most evaluations are conducted on base models and that the effect of instruction tuning on detector performance is underexplored --- precisely the gap we address.

\subsection{Hallucination Benchmarks}

TruthfulQA \cite{Lin2022TruthfulQA} introduced adversarially constructed questions designed to elicit imitative falsehoods from LLMs, and showed that larger models are not necessarily more truthful. This benchmark focuses on model-level propensity to hallucinate rather than instance-level detection.

HaluEval \cite{Li2023HaluEval} provides 35,000 hallucinated QA, dialogue, and summarization samples paired with human-verified correct answers, designed for evaluating hallucination detection methods. The QA split provides binary labels (correct vs.\ hallucinated) at the instance level, making it suitable for AUROC evaluation of detectors. We use a 1,000-question balanced sample from this split.

\subsection{RLHF Fine-Tuning and Its Effects}

RLHF fine-tuning \cite{Ouyang2022Training,Ziegler2019Fine} aligns language models with human preferences by training a reward model on human comparisons and then optimizing the LM against that reward via PPO. Constitutional AI \cite{Bai2022Training} extends this with self-critique and revision steps. Both approaches produce models that are fluent, helpful, and confident-sounding --- properties that may suppress the generative variance that sampling-based detectors rely on.

To our knowledge, no prior work has directly studied how RLHF fine-tuning affects the discriminative power of SMC-style detectors. Our work fills this gap by comparing detector performance on an RLHF-fine-tuned model against the regime assumptions made by the original SMC papers.

\subsection{NLI and Embedding Models}

We use DeBERTa \cite{He2021DeBERTa} as our NLI backbone, following prior work that established its strong performance on natural language inference tasks. For embedding-based consistency, we use Sentence-BERT \cite{Reimers2019Sentence}, which produces semantically meaningful sentence embeddings suitable for cosine similarity comparison. Both models are used in their standard pretrained configurations without further fine-tuning.
